\section{Related Work}

\subsection{Execution Traces for Code LLMs}

Recent work has scaled trace-augmented training: TracePile \citep{chen2025tracepile} pre-trained on 2.6M execution-annotated samples, while \citet{liu2023codeexecution} pre-trained models to predict execution outputs from traces.
\citet{wang2025dcsh} found trace utility to be task- and format-dependent but did not decompose individual trace components; relative to their work, we contribute a two-tier signal hierarchy, a phase transition in trace-length budgets, calibration analysis, and cross-architecture validation.
NExT \citep{ni2024next} and ET-CoT \citep{eguchi2025etcot} used trace chain-of-thought for execution reasoning, and \citet{ferrando2025execute} found compact traces comparable to full traces, consistent with our density finding.
CRUXEval \citep{gu2024cruxeval}, extended to 19 languages by CRUXEval-X \citep{xu2025cruxevalx}, established execution reasoning benchmarks, though neither decomposes which trace components contribute to \emph{classification} accuracy.
\citet{shi2023distracted} showed LLMs are easily distracted by irrelevant context, providing a complementary explanation for why verbose traces impair SFT learning.
In the prompt compression literature, LLMLingua \citep{jiang2023llmlingua} and Selective Context \citep{li2023selectivecontext} showed that pruning low-information tokens preserves downstream accuracy; TraceCoder \citep{huang2026tracecoder} and InspectCoder \citep{wang2025inspectcoder} apply selective trace extraction for automated debugging; and \citet{haque2025traces} found raw traces provide limited repair benefit.
Our density decomposition complements these methods by quantifying \emph{why} compression helps (gradient dilution explains ${\sim}$24\% of the gap; the remainder reflects inference-time interference).

\subsection{Verdict Prediction and Fault Detection}

Verdict prediction is an instance of the longstanding test oracle problem \citep{barr2015oracle}.
TRACED \citep{ding2024traced} showed execution-aware pretraining improves vulnerability detection; DeepDFA \citep{steenhoek2024deepdfa} and learned program slicing \citep{yadavally2024slicing} combined classical program analysis with deep learning for defect detection.
In patch correctness prediction, representation learning of code changes has progressed from CC2Vec \citep{hoang2020cc2vec} to LLM4PatchCorrect \citep{zhou2024llm4patch} (84.4\%) and APPT \citep{appt2024} (79.7\%); \citet{tian2020patchcorrectness} evaluated change representations for predicting repair correctness.
FuseFL \citep{widyasari2024fusefl} showed that structured fault signals such as stack traces improve localization, and \citet{foster2025mutation} leveraged mutation outcomes to guide test generation; the equivalent-mutant problem \citep{papadakis2019mutation} poses an analogous classification challenge.
Spectrum-based fault localization (SBFL) ranks program elements by suspiciousness from test spectra \citep{abreu2007accuracy,jones2002visualization,wong2014dstar}; recent LLM-based approaches such as AutoFL \citep{kang2024autofl} and AgentFL \citep{qin2024agentfl} share our insight that anomalous execution events carry the strongest diagnostic signal.
In agent-based repair, RepairAgent \citep{bouzenia2025repairagent}, Agentless \citep{xia2024agentless}, SWE-agent \citep{yang2024sweagent}, and OpenHands \citep{wang2025openhands} all incorporate execution feedback in repair workflows; in the SWE-bench ecosystem \citep{jimenez2024swebench}, these agents rely on test verdict signals to guide patch selection, yet which trace components carry that signal most efficiently remains unstudied.
LDB \citep{zhong2024ldb} applies segment-level trace verification for debugging, and SelfDebug \citep{chen2024selfdebug} and Reflexion \citep{shinn2023reflexion} iteratively refine code using execution outcomes.
None of these works isolate individual trace components or compare against rule-based baselines.

\begin{table}[t]
\centering
\caption{Cross-model verdict accuracy (\%, $N{=}1{,}190$) across four models and three architecture families. Bold: best per model. $^\ddag$2-seed mean (seed 42 training failure; Appendix~\ref{sec:tokenizer_analysis}). Multi-seed variance in Appendix~\ref{sec:full_results}.}
\label{tab:cross_model}
\small
\resizebox{\columnwidth}{!}{%
\begin{tabular}{lcccc}
\toprule
\textbf{Condition} & \textbf{Qwen-7B} & \textbf{Qwen-1.5B} & \textbf{CL-13B} & \textbf{DS-6.7B} \\
\midrule
exception\_only & \textbf{99.83} & \textbf{99.83} & \textbf{99.89} & \textbf{99.39} \\
ae\_only & 99.64 & 98.83 & 99.08 & 98.53 \\
VOTA & 99.05 & 99.54 & 98.74 & 99.25$^\ddag$ \\
loop\_only & 86.05 & 79.54 & 77.40 & 80.31 \\
no\_trace & 75.13 & 72.35 & 76.47 & 73.36 \\
\midrule
Majority baseline & \multicolumn{4}{c}{71.51} \\
\bottomrule
\end{tabular}%
}
\end{table}

\subsection{Calibration and Training Dynamics}

Neural network calibration \citep{guo2017calibration}, extended to LLMs \citep{kadavath2022language,geng2024calibration}, studies confidence--correctness alignment; \citet{spiess2025calibration} specifically addressed calibration for code models.
Calibration analysis of trace-based code models remains largely unexplored.
Just-in-time defect prediction \citep{kamei2013large,hoang2019deepjit,pornprasit2021jitline} shares our fault detection motivation but operates on code change features rather than execution traces; our work complements this line by showing that execution-time fault signals can be classified with near-perfect accuracy when properly extracted.
Curriculum learning \citep{bengio2009curriculum}, dataset cartography via training dynamics \citep{swayamdipta2020cartography}, and data selection methods \citep{sachdeva2024dataefficient,xie2023doremi} demonstrate that data quality and ordering improve convergence, consistent with our finding that compact exceptions enable convergence within 50 steps while verbose traces trap the model in oscillation across all 444 training steps.
\citet{yang2024generalization} further showed that fine-tuned LLMs generalize better on concise inputs, providing a complementary perspective on why information density governs SFT effectiveness.

